\documentclass[12pt,a4paper]{article}

\usepackage{geometry}
\geometry{margin=2.5cm}

\usepackage{longtable}
\usepackage{array}

\usepackage{hyperref}
\hypersetup{
colorlinks=true,
linkcolor=blue,
unicode=true
}

\usepackage{polyglossia}
\setmainlanguage{hebrew}
\setotherlanguage{english}

\usepackage{fontspec}
\defaultfontfeatures{Ligatures=TeX}
\setmainfont{FreeSerif}
\newfontfamily\hebrewfont[Script=Hebrew]{FreeSerif}
\newfontfamily\englishfont{FreeSerif}

\newcommand{\en}[1]{\textenglish{\LR{#1}}}

\title{%
\textbf{\Large סיכום השינויים באונטולוגיית כתבי־יד עבריים}\\[0.5em]
\normalsize מן המודל המקורי (פרופ' גילה פרבור) לגרסה המורחבת \en{(v1.4)}
}

\author{אלכסנדר גולדברג}
\date{פברואר 2026}

\begin{document}
\maketitle

\section{מטרת המסמך}
מסמך זה נכתב כדי לתאר בצורה תמציתית וברורה את ההבדלים העיקריים בין האונטולוגיה המקורית שפותחה על ידך לבין הגרסה החדשה שפותחה מאז \en{(v1.4)}, ולהסביר \textbf{מדוע כל שינוי בוצע}. 

העיקרון המנחה היה: \textbf{לשמר את הליבה התיאורטית של המודל המקורי (CIDOC-CRM + FRBRoo/LRMoo + אירועיות)} ולבצע הרחבות/התאמות רק היכן שהצטברו צרכים חדשים: משוב מומחים, מקרי מבחן מורכבים מהשטח, ואילוצי עבודה מול נתונים שמגיעים גם ממקורות חישוביים.

\section{מה נשמר מן המודל המקורי}
\begin{itemize}
\item \textbf{תשתית תקנית}: שמירה על שימוש ב-\textenglish{CIDOC-CRM} וב-\textenglish{FRBRoo/LRMoo}.
\item \textbf{גישה אירועית}: תיאור ``מה קרה לכתב היד'' לאורך הזמן (יצירה/העתקה, בעלות, משמורת, שינויים).
\item \textbf{הפרדת רבדים ביבליוגרפיים}: \textenglish{Work/Expression/Manifestation/Item} כמנגנון בסיס לתיאור תוכן מול נשא.
\item \textbf{הרחבות ייעודיות לכתבי־יד עבריים}: קולופון, סוגי כתב, חומרים, מאפיינים קודיקולוגיים/פליאוגרפיים.
\end{itemize}

\section{השינויים המרכזיים ולמה הם בוצעו}

\subsection{אינטגרציה עם מאגר הרשויות הלאומי ``מזל'' \en{(NLI Authorities)}}
\textbf{מה השתנה}: בגישה החדשה ניתן להשתמש ישירות ב-\textenglish{URI} של הספרייה הלאומית \en{(NLI)} עבור ישויות כמו אנשים/מקומות/יצירות, במקום \textenglish{URI} מקומי.

\textbf{למה}: כדי להבטיח זיהוי חד־משמעי, למנוע כפילויות, ולשפר קישוריות ל-\textenglish{VIAF/Wikidata/ORCID}. בנוסף, זה מקל על שילוב נתונים ממקורות ספרניים קיימים.

\subsection{מנגנון ודאות + מקור ייחוס \en{(Certainty \& Attribution Source)}}
\textbf{מה השתנה}: נוספו תכונות כלליות שמאפשרות לכל טענה/ייחוס לשאת:
\begin{itemize}
\item \textenglish{degree of certainty} (קטגוריאלי וגם מספרי)
\item \textenglish{attribution source} (מומחה/קטלוג/קולופון/\textenglish{AI})
\item \textenglish{attributed by} (מי או מה ביצעו את הייחוס)
\end{itemize}

\textbf{למה}: כדי לטפל במציאות שבה חלק מהנתונים הם תצפית ישירה, חלקם הכרעת חוקר, וחלקם תוצאה חישובית (למשל \textenglish{AI}). המנגנון מאפשר גם להציג למשתמשים את איכות המידע ולבצע שאילתות שמסננות לפי ודאות/מקור.

\subsection{מסגרת אפיסטמולוגית ושרשרת ראיות \en{(Epistemological Framework \& Evidence Chain)}}
\textbf{מה השתנה}: נוספה שכבה שמבחינה באופן שיטתי בין:
\begin{itemize}
\item \textbf{עובדה תצפיתית} (למשל מספר דפים, חומר)
\item \textbf{פרשנות מחקרית} (למשל ייחוס מקום/זמן על סמך פליאוגרפיה)
\item \textbf{גזירה חישובית} (למשל תוצאה ממערכת אוטומטית)
\end{itemize}
וכן נוספה אפשרות לתיעוד \textbf{שרשרת ראיות} המפרטת מקורות, צעדי היסק, שיטה, ועוצמת ראיה.

\textbf{למה}: זה פותר ערבוב נפוץ בקטלוגים בין ``נתון'' לבין ``מסקנה'', ומאפשר שקיפות מלאה במיוחד כאשר מוסיפים שכבות \textenglish{AI} או כאשר יש פרשנויות מתחרות.

\subsection{מבנה קודיקולוגי מקונן \en{(Nested Codicological Structure)}}
\textbf{מה השתנה}: בנוסף למודל הבסיסי \textenglish{BU/CU/PU}, נוספה אפשרות לקינון רקורסיבי של יחידות קודיקולוגיות (כלומר \textbf{CU בתוך CU} בעומק שרירותי), עם תיאור סוג היררכיה ועומקה.

\textbf{למה}: מקרי מבחן מהשטח (כתבי יד מורכבים, \textenglish{MTM/Sammelband}, אנתולוגיות, איגודים מחדש) לא תמיד מתיישבים עם היררכיה קשיחה של שלוש רמות בלבד. הקינון מאפשר לייצג מבנים שבהם יש ``תת־חטיבות'' בתוך יחידות פיזיות גדולות יותר, בלי לאבד מידע.

\subsection{פרדיגמה כפולה: קטלוגית מול פילולוגית \en{(Dual Paradigm Support)}}
\textbf{מה השתנה}: המודל החדש מאפשר לשמור במקביל:
\begin{itemize}
\item \textbf{תצוגה קטלוגית}: \textenglish{Work/Expression} לטובת קטלוג, קישור לקטלוגים, ותאימות לעולם ספרני.
\item \textbf{תצוגה פילולוגית}: שכבה שמדגישה מסורת טקסטואלית ועדות (\textenglish{TextTradition / Witness}) בהתאם לגישות ``פילולוגיה חדשה'' שבהן ``יצירה'' אינה תמיד ישות יציבה.
\item \textbf{גשר מפורש} בין שתי הפרדיגמות (כדי לאפשר שאילתות משולבות).
\end{itemize}

\textbf{למה}: משיקולים מחקריים, יש מתח בין מודל ``יצירה'' אחידה לבין מצב שבו כתבי־יד משקפים טקסט זורם, גרסאות, ותהליכי העתקה/עריכה. המודל הכפול מאפשר גם עבודה ספרנית וגם ניתוח פילולוגי מתקדם, בלי להכריח בחירה חד־משמעית.

\subsection{גרנולריות מלאה של מיקום בטקסט + עוגן חזותי \en{(Full Granularity \& IIIF)}}
\textbf{מה השתנה}: נוספה היררכיה של \textbf{מיקום בטקסט} מרמת דף ועד רמת תו, כולל טווחים, ובנוסף אפשרות להפנות ל-\textenglish{IIIF region} כעוגן חזותי.

\textbf{למה}: כדי לתמוך הן במחקר ידני מדויק והן בכלים אוטומטיים שמספקים תוצרים עדינים (שורה/מילה/תו), וגם כדי לקשר תיאור טקסטואלי לצילום.

\subsection{היררכיות קנוניות לייחוס כיסוי טקסט \en{(Canonical References)}}
\textbf{מה השתנה}: נוספה שכבה שמאפשרת לקשר \textenglish{Expression} לטווחים קנוניים (מקרא/משנה/תלמוד וכו') באופן מובנה.

\textbf{למה}: זה מאפשר שפה אחידה לשאילתות (``אילו כתבי־יד מכסים מסכת X דף Y?''), וגם מאפשר שילוב טוב יותר עם מערכות קיימות והיררכיות טקסטואליות.

\section{טבלת השוואה קצרה (מושגי־על ודפוסי מידול)}
\renewcommand{\arraystretch}{1.2}
\begin{longtable}{|p{4.2cm}|p{4.2cm}|p{5.2cm}|}
\hline
\textbf{במודל המקורי} & \textbf{במודל החדש \en{(v1.4)}} & \textbf{הנימוק לשינוי} \\
\hline
\endhead

שלוש רמות \textenglish{BU/CU/PU} (ללא קינון) &
\textenglish{BU/CU/PU} + \textbf{CU מקונן} &
ייצוג כתבי־יד מורכבים שלא ניתנים לתיאור בהיררכיה קשיחה. \\
\hline

אירועים ייעודיים רבים (למשל \textenglish{CensorshipEvent}) &
שימוש מקסימלי ב-\textenglish{CIDOC-CRM} לאירועים + הרחבות רק לפי צורך &
אינטרופרביליות גבוהה יותר, מניעת ``המצאה מחדש'', גמישות להרחבות עתידיות. \\
\hline

אין הבחנה פורמלית בין עובדה לפרשנות &
\textbf{סטטוס אפיסטמולוגי} + \textbf{שרשרת ראיות} &
שקיפות מדעית, טיפול בפרשנויות מתחרות, תמיכה בנתונים חישוביים. \\
\hline

אין/מעט מנגנון ודאות ממוסד &
\textbf{ודאות} + \textbf{מקור ייחוס} + \textbf{מי ייחס} &
הפיכת איכות המידע לנתון שניתן לסינון ולשאילתה; מענה לצרכי \textenglish{AI}. \\
\hline

התבססות על \textenglish{URI} מקומי נפוץ &
שימוש מועדף ב-\textbf{NLI/Mazal URI} לפי אפשרות &
מניעת כפילויות ושיפור קישוריות למאגרים חיצוניים. \\
\hline

פרדיגמה קטלוגית מובהקת \en{(Work-based)} &
\textbf{פרדיגמה כפולה} + גשר מפורש &
תמיכה גם בקטלוג וגם במחקר פילולוגי שבו ``Work'' אינו תמיד יציב. \\
\hline

מיקום כללי יחסית &
\textbf{מיקום עד רמת תו} + \textenglish{IIIF} &
תמיכה במחקר ובאוטומציה \en{(OCR/HTR/NER)}, וקישור לצילום. \\
\hline

\end{longtable}

\section{נקודות פתוחות והצעות להמשך}
\begin{itemize}
\item \textbf{איזון בין אירועים ייעודיים לבין טיפוסי אירועים}: אם נרצה להבטיח ``קריאות'' גבוהה למומחים שאינם טכניים, אפשר להוסיף תת־מחלקות ייעודיות לאירועים מסוימים (למשל צנזורה) או לחלופין להישאר עם \textenglish{E7 Activity} ולתייג באמצעות \textenglish{P2 has type}. שתי הגישות תואמות; ההבדל הוא בעיקר בשכבת הממשק/דאטה.
\item \textbf{דוגמאות מוסכמות}: לקבוע סט קטן של כתבי־יד (``מקרי מבחן'') ולהגדיר עבורם דפוסי ייצוג מוסכמים, כדי לוודא עקביות בין עבודת קטלוג, מחקר, ואוטומציה.
\end{itemize}

\section{סיכום}
הגרסה החדשה נועדה להיות \textbf{תואמת־תקנים, עשירה מספיק לכתבי־יד מורכבים, ושקופה מבחינה אפיסטמולוגית}. היא שומרת על יסודות המודל המקורי ומוסיפה בעיקר יכולות שמאפשרות: קינון מבני, עבודה בפרדיגמות מחקר שונות, ואיכות־נתונים מפורשת (ודאות/מקור/ראיות) — במיוחד כאשר נכנסים רכיבים חישוביים לתהליך.

\end{document}

